%! Author = Omar Iskandarani
%! Date = 4/19/2026
%! Affiliation = Independent Researcher, Groningen, The Netherlands
%! ORCID = 0009-0006-1686-3961
%! DOI = 10.5281/zenodo.xxx

\newcommand{\paperdoi}{10.5281/zenodo.xxx}
\newcommand{\papertitle}{A Canonical Master Equation for Swirl--String Theory\\
\large Unified Mass, Clock, Geometric Gate, and Topological Kernel}
\newcommand{\paperkeywords}{PAPER_KEYWORDS}


%======================== SST MASTER EQUATION ARTICLE ========================
\documentclass[11pt,a4paper]{article}

\usepackage[T1]{fontenc}
\usepackage{lmodern}
\usepackage{amsmath,amssymb,amsfonts,amsthm,mathtools,bm}
\usepackage{physics}
\usepackage{booktabs}
\usepackage{enumitem}
\usepackage{microtype}

% --- SST macro prelude ---
\newcommand{\swirlarrow}{\mathchoice{\mkern-2mu\scriptstyle\boldsymbol{\circlearrowleft}}{\mkern-2mu\scriptstyle\boldsymbol{\circlearrowleft}}{\mkern-2mu\scriptscriptstyle\boldsymbol{\circlearrowleft}}{\mkern-2mu\scriptscriptstyle\boldsymbol{\circlearrowleft}}}
\newcommand{\swirlarrowcw}{\mathchoice{\mkern-2mu\scriptstyle\boldsymbol{\circlearrowright}}{\mkern-2mu\scriptstyle\boldsymbol{\circlearrowright}}{\mkern-2mu\scriptscriptstyle\boldsymbol{\circlearrowright}}{\mkern-2mu\scriptscriptstyle\boldsymbol{\circlearrowright}}}
\newcommand{\vswirl}{{\mathbf{v}_{\swirlarrow}}}
\newcommand{\vswirlcw}{{\mathbf{v}_{\swirlarrowcw}}}
\newcommand{\SwirlClock}{{S_t^{\swirlarrow}}}
\newcommand{\SwirlClockcw}{{S_t^{\swirlarrowcw}}}
\newcommand{\rhoF}{\rho_{\!f}}
\newcommand{\rhoE}{\rho_{\!E}}
\newcommand{\rhoM}{\rho_{\!m}}
\providecommand{\rc}{r_c}

\newtheorem{definition}{Definition}
\newtheorem{lemma}{Lemma}
\newtheorem{theorem}{Theorem}
\newtheorem{remark}{Remark}


\usepackage{siunitx}
\usepackage[hidelinks]{hyperref}
\usepackage[a4paper,margin=1in]{geometry}
\usepackage[absolute,overlay]{textpos}

\usepackage[utf8]{inputenc}
\usepackage{pgfplots}
\pgfplotsset{compat=1.18}
\usetikzlibrary{pgfplots.groupplots}

\begin{document}
    \thispagestyle{empty}%
    \begin{center}
        \Large \bfseries \papertitle \par \vspace{1cm}
        {\Large \itshape \textbf{Omar Iskandarani}\textsuperscript{\textbf{*}} \par} \vspace{0.5cm}
        {\small Independent Researcher, Groningen, The Netherlands\par} \vspace{0.5cm}
        {\today \par}  \vspace{0.5cm}
    \end{center}
    
    \begin{abstract}
        We derive a compact master equation for Swirl--String Theory (SST) that absorbs the currently separated sectors of mass generation, geometric shielding, swirl-clock kinematics, and topological state selection into a single hierarchical formula. The resulting expression factors into four pieces: a medium scale, a geometric gate, a topological kernel, and a local clock-impedance factor. We show that the formula is dimensionally exact, reduces to the existing SST mass-functional and swirl-clock sectors in the appropriate limits, and numerically reproduces the canonical geometric gate
        \(
        \lambda_c/(\pi \rc)=4/\alpha
        \)
        using the canonical SST constants. The result is proposed as a canonical organizing equation rather than as a completed full derivation of all coefficient values.
    \end{abstract}
    Keywords: \paperkeywords


    \begin{textblock*}{1\textwidth}(0.15\textwidth,1.1\textheight)\raggedright\footnotesize\renewcommand{\arraystretch}{1.0} \noindent\rule{\textwidth}{0.4pt}\\begin{equation}0.5em]
        \textsuperscript{\textbf{*}} Email: \texttt{info@omariskandarani.com}\\  
        ORCID: \texttt{\href{https://orcid.org/0009-0006-1686-3961}{0009-0006-1686-3961}}\\
        DOI: \href{https://doi.org/\paperdoi}{\paperdoi}    
    \end{textblock*}






        

        
        \section{Purpose}
            
            Several SST papers presently encode closely related structures in different notations:
            \begin{itemize}[leftmargin=1.5em]
\item mass as localized swirl energy,
\item the Swirl-Clock as a local slowing factor,
\item the shielding gate as the ratio \(\lambda_c/(\pi \rc)\),
\item topological state dependence through crossing, ropelength, helicity, and Golden-layer factors.
            \end{itemize}
            
            The aim of this article is to replace those parallel formulations by one canonical equation.
        
        \section{Primitive definitions}
            
            We take as primitive:
            \begin{equation}
                \rhoE(x)=\frac12 \rhoF(x)\,\norm{\vswirl(x)}^2,
                \qquad
                \rhoM(x)=\frac{\rhoE(x)}{c^2},
            \end{equation}
            so that
            \begin{equation}
                \rhoM(x)=\frac12 \rhoF(x)\frac{\norm{\vswirl(x)}^2}{c^2}.
            \end{equation}
            
            The local matter-branch swirl clock is
            \begin{equation}
                \SwirlClock(x)=\sqrt{1-\frac{\norm{\vswirl(x)}^2}{c^2}}.
            \end{equation}
            
            We also define the binary geometric exposure gate
            \begin{equation}
                G(K)\in\{0,1\},
            \end{equation}
            and a dimensionless topological kernel \(\Xi_K\).
            
            \begin{definition}[Topological kernel]
                The canonical dimensionless topology kernel is taken to be
                \begin{equation}
                    \Xi_K=
                    \left[\alpha_C\,C(K)+\beta_L\,L(K)+\gamma_H\,\mathcal H(K)\right]\varphi^{-2k(K)},
                \end{equation}
                where \(C(K)\) is a crossing-type invariant, \(L(K)\) a ropelength-type invariant,
                \(\mathcal H(K)\) a helicity/Hopf-type invariant contribution, and \(k(K)\) the discrete Golden-layer index.
            \end{definition}
        
        \section{Geometric gate lemma}
            
            \begin{lemma}[Canonical geometric shielding gate]
                The geometric sector factor is
                \begin{equation}
                    \Pi_K=\left(\frac{\lambda_c}{\pi \rc}\right)^{G(K)}.
                \end{equation}
                Using the SST closure relation, this is equivalently
                \begin{equation}
                    \Pi_K=\left(\frac{4}{\alpha}\right)^{G(K)}.
                \end{equation}
            \end{lemma}
            
            \begin{proof}
                The existing SST geometric closure identifies the dimensionless shielding ratio as
                \begin{equation}
                    \frac{\lambda_c}{\pi \rc}=\frac{4}{\alpha}.
                \end{equation}
                Since \(G(K)\) is binary, the state either remains unamplified (\(G=0\)) or receives one geometric exposure factor (\(G=1\)).
            \end{proof}
        
        \section{Master equation}
            
            \begin{theorem}[Canonical SST master equation]
                For a topological state \(K\) and local position \(x\), define
                \begin{equation}
                    M_K(x)=\rhoM(x)\,\rc^3\,\Pi_K\,\Xi_K\,\SwirlClock(x)^{-2}.
                \end{equation}
                Equivalently,
                \begin{equation}
                    \boxed{
                        M_K(x)=
                        \left(\frac{\rhoF(x)\,\norm{\vswirl(x)}^2}{2c^2}\right)\rc^3
                        \left(\frac{\lambda_c}{\pi \rc}\right)^{G(K)}
                        \left[\alpha_C\,C(K)+\beta_L\,L(K)+\gamma_H\,\mathcal H(K)\right]
                        \varphi^{-2k(K)}
                        \SwirlClock(x)^{-2}
                    }
                \end{equation}
                with
                \begin{equation}
                    \SwirlClock(x)=\sqrt{1-\frac{\norm{\vswirl(x)}^2}{c^2}}.
                \end{equation}
            \end{theorem}
            
            \subsection*{Dimensional check}
                
                We verify the units of the prefactor:
                \begin{equation}
                    \left[\frac{\rhoF\,\norm{\vswirl}^2}{c^2}\rc^3\right]
                    =
                    \left[\frac{\mathrm{kg\,m^{-3}}\cdot \mathrm{m^2\,s^{-2}}}{\mathrm{m^2\,s^{-2}}}\cdot \mathrm{m^3}\right]
                    =
                    \mathrm{kg}.
                \end{equation}
                All remaining factors are dimensionless. Therefore \(M_K\) has units of mass.
        
        \section{Reduction limits}
            
            \subsection{Static spectrum limit}
                
                If the local state is approximately stationary with
                \begin{equation}
                    \rhoF(x)\to \rhoF,\qquad \norm{\vswirl(x)}\to \norm{\vswirl},\qquad \SwirlClock(x)\to 1,
                \end{equation}
                then
                \begin{equation}
                    M_K
                    =
                    \left(\frac{\rhoF\norm{\vswirl}^2}{2c^2}\right)\rc^3
                    \left(\frac{\lambda_c}{\pi \rc}\right)^{G(K)}
                    \Xi_K.
                \end{equation}
                This is the canonical static mass-spectrum form.
            
            \subsection{Geometric-mass limit}
                
                If one packages the topological kernel into a ropelength/topology factor
                \begin{equation}
                    \Xi_K \longrightarrow \widetilde{L}_{\mathrm{tot}}(K)\,\mathcal K(K),
                \end{equation}
                then
                \begin{equation}
                    M_K
                    =
                    \rhoM \rc^3
                    \left(\frac{\lambda_c}{\pi \rc}\right)^{G(K)}
                    \widetilde{L}_{\mathrm{tot}}(K)\,\mathcal K(K),
                \end{equation}
                which is the pure geometric mass-functional structure.
            
            \subsection{Clock-only limit}
                
                If the state label \(K\) is fixed and only the local swirl field varies, then
                \begin{equation}
                    M_K(x)\propto \SwirlClock(x)^{-2}.
                \end{equation}
                Thus the same local swirl field that slows the clock also strengthens the local effective mass scale.
        
        \section{Weak-swirl expansion}
            
            Since
            \begin{equation}
                \SwirlClock^{-2}(x)= \frac{1}{1-\norm{\vswirl(x)}^2/c^2},
            \end{equation}
            one has for \(\norm{\vswirl}^2/c^2\ll 1\),
            \begin{equation}
                \SwirlClock^{-2}(x)
                =
                1+\frac{\norm{\vswirl(x)}^2}{c^2}
                +\mathcal O\!\left(\frac{\norm{\vswirl}^4}{c^4}\right).
            \end{equation}
            Hence the leading correction is positive and quadratic in the local swirl speed.
        
        \section{Numerical benchmark using canonical SST constants}
            
            Use
            \begin{equation}
                \norm{\vswirl}=1.09384563\times 10^6\ \mathrm{m\,s^{-1}},
                \qquad
                \rc=1.40897017\times 10^{-15}\ \mathrm{m},
                \qquad
                \rhoF=7.0\times 10^{-7}\ \mathrm{kg\,m^{-3}}.
            \end{equation}
            
            Then
            \begin{equation}
                \rhoE=\frac12 \rhoF \norm{\vswirl}^2
                \approx 4.1877439\times 10^5\ \mathrm{J\,m^{-3}},
            \end{equation}
            and
            \begin{equation}
                \rhoM=\frac{\rhoE}{c^2}
                \approx 4.6594935\times 10^{-12}\ \mathrm{kg\,m^{-3}}.
            \end{equation}
            
            The bare mass scale is
            \begin{equation}
                M_\star=\rhoM \rc^3 \approx 1.303\times 10^{-56}\ \mathrm{kg}.
            \end{equation}
            
            The geometric gate is
            \begin{equation}
                \frac{\lambda_c}{\pi \rc}\approx 5.48144\times 10^2,
            \end{equation}
            which matches
            \begin{equation}
                \frac{4}{\alpha}\approx 5.48144\times 10^2.
            \end{equation}
            
            The Swirl-Clock at the canonical swirl speed is
            \begin{equation}
                \SwirlClock=\sqrt{1-\frac{\norm{\vswirl}^2}{c^2}}
                \approx 0.999993344,
            \end{equation}
            hence
            \begin{equation}
                \SwirlClock^{-2}\approx 1.0000133.
            \end{equation}
            
            Therefore:
            \begin{enumerate}[leftmargin=1.5em]
\item the absolute mass scale is set by \(\rhoF,\vswirl,\rc\),
\item the hierarchy must predominantly come from \(\Pi_K\Xi_K\),
\item the clock factor is a small but coherent impedance correction in the canonical regime.
            \end{enumerate}
        
        \section{Interpretation}
            
            The master equation separates SST into four logically distinct sectors:
            \begin{equation}
                \boxed{
                    \text{medium scale}
                    \times
                    \text{geometric gate}
                    \times
                    \text{topological kernel}
                    \times
                    \text{clock impedance}
                }
            \end{equation}
            
            More explicitly,
            \begin{equation}
                M_K(x)
                \sim
                (\text{localized swirl energy density})
                \times
                (\text{core volume})
                \times
                (\text{state geometry})
                \times
                (\text{topological identity})
                \times
                (\text{local time impedance}).
            \end{equation}
        
        \section{Status and assumptions}
            
            \subsection*{Derived / structurally fixed}
                \begin{equation}
                    \rhoE=\frac12 \rhoF \norm{\vswirl}^2,
                    \qquad
                    \rhoM=\rhoE/c^2,
                    \qquad
                    \SwirlClock=\sqrt{1-\norm{\vswirl}^2/c^2},
                    \qquad
                    \frac{\lambda_c}{\pi \rc}=\frac4\alpha.
                \end{equation}
            
            \subsection*{Closure-level / modeled}
                \begin{equation}
                    G(K)\in\{0,1\},
                    \qquad
                    \Xi_K=
                    \left[\alpha_C C+\beta_L L+\gamma_H \mathcal H\right]\varphi^{-2k}.
                \end{equation}
            
            \subsection*{Open}
                The exact coefficient values \( \alpha_C,\beta_L,\gamma_H \), the optimal invariant set, and whether the clock coupling is best represented by \( \SwirlClock^{-2} \), \( \SwirlClock^{-1} \), or a weaker deformation ansatz.
        
        \section{Falsifiable consequences}
            
            \begin{enumerate}[leftmargin=1.5em]
\item If two states have identical \(\Pi_K\) and \(\Xi_K\), then their rest masses must differ only through local medium variations.
\item If the geometric gate interpretation is correct, state families should cluster by binary exposure class \(G(K)\).
\item If the clock-impedance coupling is correct, strong local swirl regions should induce a correlated increase in effective inertial scale and proper-time slowing.
            \end{enumerate}
            
            A kid analogy: a particle is like a tied spinning knot in water; its mass comes from how much spinning is packed into the knot, how complicated the knot is, and how much the local spinning slows the knot’s own rhythm.
            
            \begin{thebibliography}{99}
                
                \bibitem{Batchelor1967}
                G.~K. Batchelor, 1967,
                \textit{An Introduction to Fluid Dynamics},
                Cambridge University Press,
                permalink: \url{https://doi.org/10.1017/CBO9780511800955}.
                
                \bibitem{Saffman1992}
                P.~G. Saffman, 1992,
                \textit{Vortex Dynamics},
                Cambridge University Press,
                permalink: \url{https://doi.org/10.1017/CBO9780511624063}.
                
                \bibitem{Allen1992}
                L.~Allen, M.~W. Beijersbergen, R.~J.~C. Spreeuw, and J.~P. Woerdman, 1992,
                \textit{Orbital angular momentum of light and the transformation of Laguerre-Gaussian laser modes},
                Physical Review A \textbf{45}, 8185--8189,
                DOI: \url{https://doi.org/10.1103/PhysRevA.45.8185}.
                
                \bibitem{JacobsonMattingly2004}
                T.~Jacobson and D.~Mattingly, 2004,
                \textit{Einstein-Aether gravity: a status report},
                Foundations of Physics \textbf{43},
                permalink: \url{https://doi.org/10.1007/s10701-012-9661-5}.
                
                \bibitem{GW170817}
                B.~P. Abbott et al., 2017,
                \textit{GW170817: Observation of Gravitational Waves from a Binary Neutron Star Inspiral},
                Physical Review Letters \textbf{119}, 161101,
                DOI: \url{https://doi.org/10.1103/PhysRevLett.119.161101}.
                
                \bibitem{IskandaraniCanon}
                O.~Iskandarani, 2026,
                \textit{Swirl--String Theory Canon and related SST papers},
                Independent Research / project canon.
            
            \end{thebibliography}
    



\end{document}